\chapter{Introduction}
\label{chap:introduction}
\begingroup
\justifying
\setlength{\parindent}{2em}
\setlength{\parskip}{0.55\baselineskip}

\section{Retail manipulation as a test for robot policies}

Moving a tin of beans from one shelf to the next sounds trivial. For a mobile robot it is a chain of small, unforgiving steps. The instruction names a product whose label may differ from its neighbour's only in a few words of type. The tin sits among other packages with a few centimetres to spare. The base has to stop where the arm can reach, and a stop that is a few centimetres short turns a grasp into a miss. Refrigerated fixtures add a different kind of difficulty: a door has to be found, contacted at the right place and moved through an arc while the base stays out of the way. Each of these problems is familiar on its own. A store puts them in the same aisle.

That combination makes retail an attractive setting for evaluating vision-language-action (VLA) policies, the family of models that map camera images and a written instruction directly to robot actions. Several such models are now public, together with recipes for fine-tuning them on new robots. Whether a released, fine-tuned checkpoint does what its authors report when someone else runs it is a separate question, and one that benchmark papers rarely answer for each other.

RoboBenchMart~\cite{soshin2025robobenchmart} is a recent benchmark for exactly this setting. Built on ManiSkill3~\cite{taomaniskill3}, it generates dark-store layouts procedurally, stocks them with a library of retail products, and defines pick-and-place and door tasks for a Fetch mobile manipulator. Its authors generated motion-planned demonstrations, fine-tuned four VLA policies on them, evaluated the policies under four increasingly demanding conditions, and released the code, the scenes and some of the fine-tuned checkpoints. That release is the starting point of this dissertation.

\section{Problem statement}

A benchmark result is only as useful as it is reproducible. In robot learning the gap between ``the code is public'' and ``the result can be reproduced'' is wide. A result depends on the policy weights, on the version of the simulator and the environment code, on how the policy is served and how its actions are applied, and on details of the evaluation protocol that a paper may summarise in a sentence. A change in any of these can change the outcome, and the changes are not always visible from the outside.

This dissertation asks a narrow version of that question for RoboBenchMart. If I take two of the released checkpoints, Octo and $\pi_{0.5}$, and run them through the benchmark's atomic evaluation protocol using only what is publicly available today, what results do I get? How do they compare with the published ones? And what can and cannot be concluded from the comparison? I did not train any model. I did not re-run the upstream study. I ran the released artefacts as released, recorded what happened, and examined the protocol closely enough to say what the numbers mean.

\section{Research questions}

\begin{description}[leftmargin=2.8em,style=nextline]
    \item[\textbf{RQ1}] Can the 42 atomic evaluation configurations of RoboBenchMart be run to completion using only the public code, scenes and checkpoints? What has to be repaired, and what evidence does the run leave behind?
    \item[\textbf{RQ2}] How many episodes does each checkpoint complete in each of the five task families under each condition, and which differences between conditions and between models can be resolved with 30 episodes per configuration?
    \item[\textbf{RQ3}] How far are these results from the upstream report, and how do properties of the protocol and differences between software versions limit the conclusions that can be drawn?
\end{description}

\section{Scope}

The study covers two policies, Octo~\cite{octo_2023} and $\pi_{0.5}$~\cite{intelligence2025pi05visionlanguageactionmodelopenworld}, as fine-tuned and released by the RoboBenchMart authors. Each was evaluated on the 42 atomic configurations defined in the benchmark's evaluation script: 12 under each of the Seeds, Pose and Layout conditions and 6 under Pairing, with 30 episodes per configuration. That is 1,260 episodes per model.

Several things are outside the scope, and it is worth listing them because the original benchmark covers them. I evaluated neither $\pi_0$ nor SmolVLA, both of which appear in the upstream paper. I ran no composite tasks. I did not annotate why episodes failed, so the upstream failure taxonomy is described but not applied. I did not fine-tune or train anything, and the training conditions of the checkpoints are known to me only from the upstream paper. All experiments are in simulation.

One further boundary turned out to matter more than I expected. The code and scenes I evaluated on are newer than the checkpoints. The environment has changed in several ways since the checkpoints were published, and I did not attempt to reconstruct the version the checkpoints were trained against. The results therefore describe the released checkpoints on the current release of the benchmark, and Chapter~\ref{chap:results} discusses what that implies.

\section{Contributions}

The contributions are modest and practical.

The main one is an independent, fully documented run of the RoboBenchMart atomic protocol for two released checkpoints: 2,520 episodes, each linked to a seed, a log line, a video and a SHA-256 hash, together with the revisions of every component involved.

Getting the protocol to run required three repairs, which I made and documented. The floor environments could not be constructed because an upstream change had removed an import still in use. The evaluation client read the instruction from a wrapped environment. And the Octo server shared one observation history across all connections, which would have mixed episodes once clients ran in parallel.

The results themselves are the third contribution. Neither checkpoint completed any of 1,800 pick-and-place episodes; $\pi_{0.5}$ closed doors in about a fifth of its attempts. This is far below the upstream figures, by a margin that sampling noise cannot account for.

Finally, I read the evaluation code closely enough to correct several statements about the protocol, including one I had made myself in an earlier draft: the evaluation client locks the robot's head, not its torso. I also identify the version changes between the checkpoints and the evaluated environment that are most likely to matter, and I propose a version-aligned re-evaluation as the direct test.

\section{Structure}

Chapter~\ref{chap:literature} reviews work on robot benchmarks, retail datasets, procedural environments, demonstration generation and VLA policies, and on the reproducibility of evaluation results. Chapter~\ref{chap:system} describes the RoboBenchMart benchmark, all of which is upstream work. Chapter~\ref{chap:methodology} describes how I evaluated the checkpoints, and Chapter~\ref{chap:results} reports and discusses the results. Chapter~\ref{chap:conclusion} draws conclusions and sets out the next experiments. The appendix lists every configuration's result, the component revisions and the records kept.

\endgroup
